\section{Literatür verisiyle yazılım uygulamaları}
\label{sec:spss-b07}

\textbf{Amaç ve veri.} \texttt{b07.csv} üzerinden ortalama ve oran için
nokta tahmini üretin \citep{cortez2008data}. \texttt{pass10}, yıl sonu notu
10 veya üzerinde ise 1'dir.

\textbf{Ortak veri.} Doğrulanan B07 uygulamasında üç yazılım da
\nolinkurl{companion/spss/csv/b07.csv} dosyasını kullanmıştır. Dosya
395 satır ve 10 sütun içerir. B06'nın seçim göstergeleri ve ağırlığı bu
uygulamada kullanılmaz; 395 kaydın tamamı analiz edilir. Bu karşılaştırma
Excel içe aktarmanın doğrulandığı anlamına gelmez. Kodları
\texttt{companion} klasörünü içeren paket kökünden çalıştırın;
veri sözlüğü ve hazırlık için bkz. sayfa~\pageref{sec:spss-guide}.

\subsection{IBM SPSS uygulaması}
\textbf{Başlamadan önce.} Aşağıdaki komut dosyası B07 CSV verisini
açar ve özetler. Menü yolunu kullanacaksanız önce veri açma ve tanımlama
komutlarını \texttt{EXECUTE} satırı dahil çalıştırın. Filtre, ağırlık ve
dosya bölme kapalı olmalıdır. Komut yolundaki \texttt{AGGREGATE}, etkin
veriyi tek özet satırıyla değiştirir; bu özeti ham veri dosyasının üzerine
kaydetmeyin. Analizi tekrarlamak için kodu veri açma adımından başlatın.

\begin{spssadimlar}
\spssadim{\emph{Analyze $\to$ Descriptive Statistics $\to$ Frequencies} yolunda \texttt{pass10} seçin ve tabloyu alın. 1 kategorisinin oranını belirleyin.}
\spssadim{\emph{Data $\to$ Aggregate} yolunu açın. \emph{Break Variable(s)} alanını boş bırakın; böylece dosyanın tamamı tek grup olarak özetlenir.}
\spssadim{\texttt{g3} için \emph{Mean} ve \emph{Standard deviation}, \texttt{pass10} için \emph{Mean} işlevlerini seçin. Özetlere sırasıyla \texttt{ort}, \texttt{ss}, \texttt{oran} adlarını verin; grup gözlem sayısını \texttt{n} adıyla ekleyin.}
\spssadim{Yalnız özet değişkenleri içeren yeni bir veri kümesi oluşturmayı seçin. \emph{OK} sonrasında o veri kümesini etkinleştirin; bir satır bütün dosyanın özetidir.}
\spssadim{\emph{Transform $\to$ Compute Variable} yolunda sırayla \texttt{sh\_ort=ss/SQRT(n)} ve \texttt{sh\_oran=SQRT(oran*(1-oran)/n)} hesaplarını yapın. Sol tarafı hedef değişken, sağ tarafı sayısal ifade alanına girin.}
\spssadim{\emph{Data View} içinde özetleri inceleyin. Bir çıktı listesi için aşağıdaki tam komut dosyasını baştan çalıştırın; bu yol ham veriyi yeniden açar ve özet veriyi etkin dosyanın yerine koyar.}
\end{spssadimlar}

{\raggedright
\noindent\textbf{Komutların görevi.} \texttt{AGGREGATE}, boş \texttt{BREAK} ile bütün dosyayı tek özet satırına indirir. \texttt{COMPUTE} standart hataları hesaplar; \texttt{FORMATS} görünümü, \texttt{LIST} çıktı listesini düzenler.\par}

\textbf{Uygulama kodu:} \texttt{b07}. Doğrulanan veri dosyası
\texttt{companion/spss} altındaki \texttt{csv/b07.csv} dosyasıdır.

\begin{lstlisting}[style=prismSPSS]
SET DECIMAL=DOT.
GET DATA /TYPE=TXT
 /FILE='companion/spss/csv/b07.csv'
 /ENCODING='UTF8'
 /ARRANGEMENT=DELIMITED /DELCASE=LINE
 /DELIMITERS="," /QUALIFIER='"'
 /FIRSTCASE=2
 /VARIABLES=id F8.0 school F8.0 sex F8.0 age F8.0
 studytime F8.0 absences F8.0 g1 F8.0 g2 F8.0
 g3 F8.0 pass10 F8.0.
FILTER OFF.
WEIGHT OFF.
SPLIT FILE OFF.
VARIABLE LABELS
 g3 'Yil sonu matematik notu'
 pass10 'Yil sonu notu 10 veya uzerinde'.
VALUE LABELS pass10
 0 '10 puanin altinda' 1 '10 puan ve uzerinde'.
VARIABLE LEVEL
 id school sex pass10 (NOMINAL)
 studytime (ORDINAL)
 age absences g1 g2 g3 (SCALE).
EXECUTE.
FREQUENCIES VARIABLES=pass10.
AGGREGATE /OUTFILE=*
 /BREAK= /n=N /ort=MEAN(g3) /ss=SD(g3)
 /oran=MEAN(pass10).
COMPUTE sh_ort=ss/SQRT(n).
COMPUTE sh_oran=SQRT(oran*(1-oran)/n).
FORMATS n (F8.0).
FORMATS ort ss oran sh_ort sh_oran (F14.9).
LIST VARIABLES=n ort ss sh_ort oran sh_oran.
\end{lstlisting}

\begin{outputguidebox}
\textbf{Paylaşılan SPSS çıktısıyla doğrulanan değerler.}
\texttt{LIST} çıktısında $n=\SPContextN$, ortalama \SPContextMean,
$s=\SPContextSD$, ortalamanın standart hatası \SPContextSE\ bulunmuştur.
Eşiği karşılayan \SPContextPass\ öğrenci vardır; oran \SPContextPHat,
bağımsız Bernoulli modeli altındaki oran standart hatası \SPContextPSE'dır.
Ortalama ve oran farklı tahmin hedefleridir; örneklem oranı evren oranının
kesin değeri değildir. \texttt{AGGREGATE} veriyi tek özet satırına indirir;
çıktıdaki bir satır yalnız bir öğrencinin ölçüldüğü anlamına gelmez.
\end{outputguidebox}


\par\noindent\begin{minipage}{\linewidth}
\subsection{R uygulaması}
\label{sec:r-gercek-b07}

İkili \texttt{pass10} değişkeninin ortalaması oran tahminidir. Ortalama ve oran standart hataları farklı formüllerle hesaplanır.

\begin{lstlisting}[style=prismR]
veri <- read.csv("companion/spss/csv/b07.csv")
stopifnot(nrow(veri) == 395, ncol(veri) == 10,
          !anyNA(veri),
          all(veri$pass10 == as.integer(veri$g3 >= 10)))
gozlem_sayisi <- nrow(veri)
ortalama <- mean(veri$g3)
standart_sapma <- sd(veri$g3)
oran <- mean(veri$pass10)
frekans <- table(veri$pass10)
print(frekans); print(100 * prop.table(frekans))
print(c(n=gozlem_sayisi, ort=ortalama, ss=standart_sapma,
  sh_ort=standart_sapma/sqrt(gozlem_sayisi), oran=oran,
  sh_oran=sqrt(oran*(1-oran)/gozlem_sayisi)))
print(sum(veri$g3 == 0))
\end{lstlisting}

\textbf{Çıktıyı okuma.} Paylaşılan R çıktısı 395 satır, 10 sütun ve
her sütunda sıfır eksik değer gösterir. \texttt{pass10} ile
\texttt{g3 >= 10} koşulu bütün kayıtlarda uyumludur. \texttt{ort} ve
\texttt{oran} nokta tahminleri; \texttt{sh\_ort} ve \texttt{sh\_oran}
ise bunların kullanılan model altındaki standart hatalarıdır.
Sıfır notlu 38 kayıt analizde tutulmuştur.
\end{minipage}\par

\par\noindent\begin{minipage}{\linewidth}
\subsection{Python uygulaması}
\label{sec:python-b07}

\texttt{ddof=1} örneklem standart sapmasını seçer; \texttt{sh\_ort} ile \texttt{sh\_oran} ayrı belirsizlik ölçüleridir.

\begin{lstlisting}[style=prismPythonApp]
import pandas as pd
import numpy as np

veri = pd.read_csv("companion/spss/csv/b07.csv")
assert veri.shape == (395, 10)
assert not veri.isna().any().any()
assert veri.pass10.eq((veri.g3 >= 10).astype(int)).all()
gozlem_sayisi = len(veri)
ortalama = veri.g3.mean()
standart_sapma = veri.g3.std(ddof=1)
oran = veri.pass10.mean()
frekans = veri.pass10.value_counts().sort_index()
print(frekans); print(100 * frekans / gozlem_sayisi)
print(pd.Series({
    "n": gozlem_sayisi, "ort": ortalama, "ss": standart_sapma,
    "sh_ort": standart_sapma / np.sqrt(gozlem_sayisi),
    "oran": oran,
    "sh_oran": np.sqrt(oran*(1-oran)/gozlem_sayisi)
}))
print(int((veri.g3 == 0).sum()))
\end{lstlisting}

\textbf{Çıktıyı okuma.} Paylaşılan Python çıktısında da 395 satır,
10 sütun, sıfır eksik değer ve eşik kodlaması için tam uyum vardır.
Altı özet değer R ile eşleşir; 130 ve 265 kişilik frekanslar ile
38 sıfır notlu kayıt da doğrulanmıştır. Bu uygulama bir hipotez testi,
$p$ değeri veya güven aralığı üretmez. R ve Python ana veriyi korur.
\end{minipage}\par

\subsection{Sonuçların karşılaştırılması ve ortak rapor}
12 Eylül 2026'da paylaşılan SPSS tabloları ile R ve Python metin
çıktıları karşılaştırılmıştır. Gözlem sayısı, ortalama, standart sapma,
ortalamanın standart hatası, oran ve oranın standart hatasından oluşan
altı değer R ve Python'da gösterilen 13 ondalık basamakta eşleşmiştir.
SPSS listesindeki ondalıklı değerler dokuz basamakta aynı sonuçlarla
uyumludur. Frekanslar aynıdır; yüzdeler gösterilen yuvarlama
hassasiyetlerinde uyumludur. Proje CSV'sinden yapılan hesaplar da
bu sonuçları destekler. Aşağıdaki tablolar çıktıların kitap düzenine
uyarlanmış özetidir; ekran görüntüsü değildir.

\begin{table}[H]
\centering
\caption{B07'de yıl sonu notunun 10 puan eşiğine göre dağılımı}
\label{tab:b07-dogrulanmis-frekans}
\begin{tabular}{@{}lrr@{}}
\toprule
Kategori & Frekans & Yüzde \\
\midrule
10 puanın altında & 130 & 32,911392 \\
10 puan ve üzerinde & 265 & 67,088608 \\
Toplam & 395 & 100,000000 \\
\bottomrule
\end{tabular}
\par\smallskip
{\footnotesize Yüzdeler altı ondalık basamağa yuvarlanmıştır.
Eksik değer olmadığından SPSS'in \emph{Percent} ve \emph{Valid Percent}
sütunları eşittir; yüzdelerin paydası 395 kayıttır.\par}
\end{table}

\begin{table}[H]
\centering
\caption{B07'nin doğrulanmış nokta tahminleri ve standart hataları}
\label{tab:b07-dogrulanmis-tahmin}
\begin{tabular}{@{}lr@{}}
\toprule
Ölçü & Değer \\
\midrule
Gözlem sayısı & 395 \\
Ortalama not & 10,415190 \\
Notların standart sapması & 4,581443 \\
Ortalamanın standart hatası & 0,230517 \\
10 puan eşiğini karşılayanların oranı & 0,670886 \\
Oranın standart hatası & 0,023643 \\
\bottomrule
\end{tabular}
\par\smallskip
{\footnotesize Ondalıklı değerler altı basamağa yuvarlanmıştır.
Standart sapma 394 böleniyle hesaplanmıştır. Standart hatalar
$s/\sqrt{395}$ ve $\sqrt{\hat p(1-\hat p)/395}$ formüllerine dayanır;
ara hesaplarda yuvarlanmamış değerler kullanılmıştır.\par}
\end{table}

\Needspace{15\baselineskip}
\begin{outputguidebox}
Tablo~\ref{tab:b07-dogrulanmis-frekans} içindeki 265 kişi,
$\hat p=265/395$ oranını verir. Tablo~\ref{tab:b07-dogrulanmis-tahmin}
içindeki 4,581443, bireysel notların yayılımını; 0,230517 ise kullanılan
bağımsız gözlemler modeli altında ortalama tahmininin standart hatasını
gösterir. Oranın standart hatası 0,023643, yaklaşık 2,3643 yüzde puanıdır;
\%0,023643 olarak okunmaz.

Bu standart hatalar öğretim amacıyla kullanılan modelin hesaplarıdır;
kaynak araştırmanın örnekleme tasarımına uygun belirsizlik hesabının
veya daha geniş bir evrene genellemenin doğrulandığı anlamına gelmez.
Yalnız eldeki 395 kaydın betimlenmesi hedefleniyorsa ortalama ve oran
bu dosyanın özetleridir. Üç yazılımın uyumu temsil gücünü kanıtlamaz.

SPSS'in \texttt{LIST} çıktısındaki bir satır, 395 öğrencinin toplu
özetidir. Ham dosyada 38 sıfır not vardır; bunlar eksik değer sayılmamış,
hem ortalama hesabında hem de eşik altı kategorisinde korunmuştur.
\end{outputguidebox}

\Needspace{9\baselineskip}
\textbf{Raporlama örneği (doğrulanmış çıktılarla).}
``395 kayıtta ortalama yıl sonu notu 10,415190, standart sapması
4,581443 bulunmuştur. Notu 10 veya üzerinde olan 265 kişinin oranı
0,670886'dır (\%67,088608). Bağımsız gözlemler modeli altında
ortalamanın standart hatası 0,230517, oranın standart hatası 0,023643
olarak hesaplanmıştır. Sıfır notlu 38 kayıt analizde tutulmuştur.
Sonuçlar kaynak dosyayı özetlemektedir; daha geniş bir öğrenci evrenine
genellenebilirlik iddiası taşımamaktadır.''


\textbf{Görev.} Nokta tahminlerinin yanına hedef evren tanımını ekleyin.
Sıfır notları gerekçesiz biçimde eksik saymanın analiz paydasını ve
tahmin hedefini nasıl değiştireceğini açıklayın. Yalnız \texttt{g3}
değişkenini değiştirmek ile aynı kayıtları \texttt{pass10} hesabından da
çıkarmak arasındaki farkı belirtin; 99 yazmanın tek başına eksik değer
tanımlamak olmadığını unutmayın. SPSS'te sonraki uygulamaya geçmeden
önce tek satırlık özet yerine gerekli ham veri dosyasını açın.